%=============================================================
%  Harmonic-Drive Digital Twin — Package Anatomy
%  Compile: latexmk -xelatex package_anatomy.tex
%=============================================================
\documentclass[11pt,a4paper]{article}
\usepackage{schaeffler}          % Schaeffler branding (Liberation Sans, green)
\usepackage{fancyvrb}

% Box-drawing glyphs in the directory trees need a mono font that has them.
\setmonofont{DejaVu Sans Mono}[Scale=0.86]

%---- metadata ------------------------------------------------
\hypersetup{
  pdftitle={Harmonic-Drive Digital Twin — Package Anatomy},
  pdfauthor={Lu Zhu},
  pdfsubject={What is inside the released package, piece by piece},
}

\newcommand{\pkg}[1]{\texttt{#1}}

\begin{document}

%==============================================================
\schaefflertitlepage
  {Harmonic-Drive Digital Twin}
  {Package Anatomy}
  {What is inside the release · piece by piece · in plain language}
  {Lu Zhu}
  {2026}
  {Rev 1.0 — Internal}

\tableofcontents
\clearpage

%==============================================================
\section{Read this first}

This document explains \textbf{what you are shipping} when you release
\pkg{harmonic-drive-twin}, piece by piece, in plain language.

It is written for two readers: \textbf{you}, deciding what belongs in the
release; and \textbf{a stranger} who has just cloned the repository and wants
to know what they are looking at.

\begin{keyresult}[The whole package in one paragraph]
The package is a \textbf{physics calculator for a strain-wave gearbox}. You
hand it an operating point --- ``I need 30\,N\,m out, the motor is spinning at
2000\,rpm, it is 20\,°C in there'' --- and it hands back the motor torque you
actually need, the efficiency, where every watt of loss went, and how much the
gearbox twists under that load. Everything else in the repository exists to
support, prove, or demonstrate that one calculation.
\end{keyresult}

Its companion, the \emph{User Guide}, tells a reader \textbf{how to drive}.
This document tells them \textbf{what is under the hood}.

%==============================================================
\section{The map}

\begin{Verbatim}[fontsize=\footnotesize,xleftmargin=6pt]
harmonic-drive-twin/
│
├── pyproject.toml       ← the shipping label (name, version, deps)
├── LICENSE              ← MIT — permission to use it
├── README.md            ← the shop window (what it does, quick start)
├── .gitignore           ← what never leaves your machine
│
├── harmonic_drive_twin/ ← THE PRODUCT — all of the model is here
│   ├── __init__.py          the front door
│   ├── twin.py              the engine
│   ├── params.py            the dials
│   ├── catalogue.py         the reality check
│   ├── py.typed             flag saying "type hints included"
│   └── flex_common/         the physics foundation (7 modules)
│       ├── baseline.py   ellipse.py    mesh.py     ring.py
│       └── friction.py   profile.py    teeth.py
│
├── examples/            ← on-ramps: three runnable scripts
│   ├── basic_usage.py
│   ├── isaacsim6_actuator.py
│   └── isaacsim6_pendulum.py
│
├── tests/               ← the proof it still works
│   └── test_catalogue_validation.py
│
└── docs/                ← the manual
    └── user_guide/
        ├── user_guide.pdf        10-page branded guide
        ├── package_anatomy.pdf   this document
        └── assets/               logos + pendulum result figure
\end{Verbatim}

Total: \textbf{2\,947 lines of Python} in the model itself (3\,633 counting
examples and tests). Small --- and that is a feature. A reviewer can read the
whole model in an afternoon.

%==============================================================
\section{The four layers}

The package is built in layers. Each layer knows only about the ones below it.
This is why it stays understandable.

\begin{Verbatim}[fontsize=\small,xleftmargin=20pt]
    ┌────────────────────────────────────────────┐
 4  │  __init__.py    the front door             │  ← users import this
    ├────────────────────────────────────────────┤
 3  │  twin.py        the engine                 │  ← does the work
    ├────────────────────────────────────────────┤
 2  │  params.py      the dials                  │  ← what you may change
    ├────────────────────────────────────────────┤
 1  │  flex_common/   the physics foundation     │  ← the equations
    └────────────────────────────────────────────┘

  catalogue.py sits alongside 2–3 and asks: "does this match reality?"
\end{Verbatim}

%--------------------------------------------------------------
\subsection{Layer 1 --- \pkg{flex\_common/} : the physics foundation}

\emph{Roughly 1\,300 lines across seven small modules. Each is one piece of
textbook mechanics.}

These are the equations describing \emph{how a harmonic drive deforms and
meshes}. They began as the shared library of the multi-topic FEA study and were
bundled into the package so it stands alone --- a user does \textbf{not} need
the research repository to run the model.

\renewcommand{\arraystretch}{1.25}
\begin{tabularx}{\textwidth}{@{}l r Y@{}}
\rowcolor{SchGreen}
\whitebf{Module} & \whitebf{Lines} & \whitebf{In one sentence} \\
\pkg{baseline.py} & 79 & \textbf{The reference part.} One agreed set of numbers
--- $R=20$\,mm, wall 0.4\,mm, ovalization 0.30\,mm, 200/202 teeth --- so every
other module is talking about the same gearbox. \\
\rowcolor{SchGreenLight}
\pkg{ring.py} & 170 & \textbf{The flexspline as a thin ring.} Squash a thin
steel ring into an ellipse; this gives the bending moment, shear and contact
pressure in closed form. Handles both ``tied to the cam'' and ``cam can only
push outward''. \\
\pkg{ellipse.py} & 111 & \textbf{A real ellipse is not a pure $\cos2\theta$.}
Decomposes the true elliptical shape into harmonics and adds the ring response
to each. The correction term for \pkg{ring.py}. \\
\rowcolor{SchGreenLight}
\pkg{mesh.py} & 305 & \textbf{Which teeth are touching, right now.} Given the
wave-generator angle it says which flexspline teeth line up with which
circular-spline spaces, and how each tooth moves: radial breathing, tangential
slide, tilt. The heart of the kinematics. \\
\pkg{profile.py} & 433 & \textbf{What shape the teeth must be.} Synthesises the
circular-spline flank as the conjugate envelope of the flexspline tooth swept
through the mesh, and computes the flank sliding speed --- which feeds friction
and wear. \\
\rowcolor{SchGreenLight}
\pkg{friction.py} & 74 & \textbf{Why it takes torque to spin an unloaded
drive.} Without friction the no-load torque is exactly zero; this shows why,
and what friction adds. \\
\pkg{teeth.py} & 184 & \textbf{Tooth-root stress.} Lewis cantilever bending
plus Dolan--Broghamer fillet concentration --- the fatigue surrogate. \\
\end{tabularx}

\vspace{6pt}
\noindent\emph{Gentle note:} a user of the twin never needs to open these. They
are the foundation under the floorboards. But they are what makes the model
\emph{physics-based} rather than a curve fit --- and a skeptical reviewer will
go straight here.

%--------------------------------------------------------------
\subsection{Layer 2 --- \pkg{params.py} : the dials}

\emph{425 lines. Four dataclasses. This is ``everything you are allowed to
change''.}

Instead of magic numbers scattered through the code, every adjustable quantity
lives in one of four labelled boxes.

\begin{tabularx}{\textwidth}{@{}l Y Y@{}}
\rowcolor{SchGreen}
\whitebf{Dataclass} & \whitebf{What it holds} & \whitebf{Example knobs} \\
\pkg{FrictionParams} & How things rub & boundary friction coefficient, Stribeck
transition speed, wave-generator bearing coefficient, grease churn factor \\
\rowcolor{SchGreenLight}
\pkg{ThermalParams} & How temperature changes things & oil viscosity versus
temperature, thermal expansion of clearances \\
\pkg{ToleranceParams} & Manufacturing imperfection & backlash, tooth-to-tooth
pitch scatter (used for Monte-Carlo runs) \\
\rowcolor{SchGreenLight}
\pkg{StiffnessChain} & How the drive twists & the series chain: cup +
wave-generator bearing + mesh + output bearing \\
\end{tabularx}

\vspace{6pt}
\noindent\textbf{Why this matters for release:} a user adapting the twin to
\emph{their} gearbox edits parameters here --- never the equations. That is the
contract. Units are stated once, at the top of the file, and hold everywhere:
\textbf{mm, N, N\,mm, rad/s, °C}.

%--------------------------------------------------------------
\subsection{Layer 3 --- \pkg{twin.py} : the engine}

\emph{787 lines. One class, \pkg{HarmonicDriveTwin}, plus three contact
helpers.}

This is the actual model. It takes the equations from Layer 1, the numbers from
Layer 2, and answers questions. The main method is \pkg{solve()}:

\begin{Verbatim}[fontsize=\small,xleftmargin=10pt]
         you give it                     it gives you back
   ┌──────────────────────┐    ┌─────────────────────────────────┐
   │  T_out    torque     │    │  T_in       motor torque needed │
   │  omega_in speed      │ ─▶ │  eta        efficiency          │
   │  T        temp °C    │    │  T_teeth_in ┐                   │
   └──────────────────────┘    │  T_wg_in    ├ where the loss    │
                               │  T_churn    ┘ actually went     │
                               │  dpsi_total how much it twists  │
                               │  K_total    torsional stiffness │
                               └─────────────────────────────────┘
\end{Verbatim}

The loss split into \textbf{three channels} is the model's central idea:

\begin{enumerate}
  \item \pkg{T\_teeth\_in} --- tooth flanks rubbing. Boundary-lubricated, not
        oil-film.
  \item \pkg{T\_wg\_in} --- the wave-generator bearing rolling.
  \item \pkg{T\_churn} --- grease churn and seal drag. Load-\emph{independent},
        and at rated torque it is the \textbf{biggest} of the three.
\end{enumerate}

\begin{keyresult}[Why the third channel matters]
The load-free churn/seal channel is the non-obvious finding the study produced,
and it is why catalogue efficiency sits at 67--78\,\% rather than the
$\approx$\,90\,\% a tooth-friction-only model predicts.
\end{keyresult}

Beyond \pkg{solve()}, the class also answers:

\begin{tabularx}{\textwidth}{@{}l Y@{}}
\rowcolor{SchGreen}
\whitebf{Method} & \whitebf{Question it answers} \\
\pkg{efficiency\_curve()} & How does $\eta$ vary with torque? \\
\rowcolor{SchGreenLight}
\pkg{efficiency\_map()} & 2-D $\eta$ map over torque $\times$ speed \\
\pkg{hysteresis()} / \pkg{hysteresis\_dahl()} & What does the torque--windup
loop look like? (lost motion) \\
\rowcolor{SchGreenLight}
\pkg{breakaway\_input\_torque()} & What torque starts it from cold rest? \\
\pkg{backdrive\_threshold()} & Can the load push the motor backwards? \\
\rowcolor{SchGreenLight}
\pkg{monte\_carlo()} & What is the scatter if tooth clearances vary? \\
\pkg{warmup()} & Thermal transient and grease life \\
\rowcolor{SchGreenLight}
\pkg{wear\_life()} & Archard flank wear over a lifetime \\
\pkg{hertz\_flank\_summary()} & Contact pressure and oil film at the
hardest-loaded tooth \\
\end{tabularx}

\vspace{6pt}
\noindent Plus three standalone helpers anyone can use:
\pkg{hertz\_flank()},\newline
\pkg{ehl\_lambda()} and \pkg{wg\_bearing\_stiffness()}.

%--------------------------------------------------------------
\subsection{Layer 4 --- \pkg{\_\_init\_\_.py} : the front door}

\emph{129 lines, almost all of it documentation.} Its job is \emph{curation}:

\begin{enumerate}
  \item \textbf{Decides what is public.} The \pkg{\_\_all\_\_} list is the
        promise: these 13 names are the supported API. Everything else is
        internal and may change.
  \item \textbf{Provides \pkg{solve\_point()}} --- a one-line functional
        wrapper so a casual user never has to instantiate a class:
        \pkg{solve\_point(T\_out\_Nm=30.0, speed\_rpm=2000)} $\rightarrow$
        $\eta = 77.9$\,\%.
  \item \textbf{Carries the quick-start docstring}, written as doctests --- so
        the examples in the documentation are \emph{executable} and cannot
        silently go stale.
\end{enumerate}

%==============================================================
\section{The reality check}

This is what separates a digital twin from a plausible-looking spreadsheet.

\subsection{\pkg{catalogue.py} (250 lines)}

Holds the digitised \textbf{Schaeffler TPI 275} table (series RT1-H-CS,
$i = 100$, $+20$\,°C) for four frame sizes:

\begin{tabularx}{\textwidth}{@{}l Y r r r r@{}}
\rowcolor{SchGreen}
\whitebf{Symbol} & \whitebf{Quantity} & \whitebf{14} & \whitebf{17} &
\whitebf{20} & \whitebf{25} \\
$R$ & pitch radius [mm] & 17.78 & 21.59 & 25.40 & 31.75 \\
\rowcolor{SchGreenLight}
$T_N$ & rated torque [N\,m] & 10.0 & 31.0 & 52.0 & 87.0 \\
$\eta$ & efficiency [\%] & 67.0 & 77.0 & 77.0 & 77.0 \\
\rowcolor{SchGreenLight}
$T_\mathrm{nlrt}$ & no-load running torque [mN\,m] & 35 & 51 & 105 & 195 \\
$T_\mathrm{nlst}$ & starting torque [mN\,m] & 21 & 29 & 37 & 69 \\
\rowcolor{SchGreenLight}
$T_\mathrm{bt}$ & back-driving torque [N\,m] & 2.21 & 3.06 & 3.89 & 7.26 \\
$K_1$ & torsional stiffness, low [N\,m/rad] & 4\,700 & 10\,000 & 16\,000 &
31\,000 \\
\rowcolor{SchGreenLight}
$K_3$ & torsional stiffness, high [N\,m/rad] & 7\,100 & 16\,000 & 29\,000 &
57\,000 \\
\end{tabularx}

\vspace{8pt}
\noindent And three functions: \pkg{twin\_for\_size(14|17|20|25)} returns a twin
pre-scaled to that catalogue size; \pkg{calibrate()} derives the three global
calibration dials from the table; \pkg{validate()} runs everything and reports
how close the model lands.

\begin{keyresult}[The headline claim of the whole package]
Five independent catalogue quantities, across four sizes, matched with only
\textbf{three global calibration dials} and \textbf{zero per-size fitting}.
Torsional stiffness uses no fitting at all --- it comes from the FE-anchored
physical chain.
\end{keyresult}

\subsection{\pkg{tests/test\_catalogue\_validation.py} (183 lines)}

Four tests that keep that claim honest:

\begin{tabularx}{\textwidth}{@{}l Y@{}}
\rowcolor{SchGreen}
\whitebf{Test} & \whitebf{Guards against} \\
\pkg{test\_calibration\_reproducible()} & calibration drifting when code
changes \\
\rowcolor{SchGreenLight}
\pkg{test\_solve\_point\_size17()} & the size-17 number (75.8\,\% versus
catalogue 77\,$\pm$\,3\,\%) moving \\
\pkg{test\_solve\_point\_default()} & the baseline number (77.9\,\%) moving \\
\rowcolor{SchGreenLight}
\pkg{test\_twin\_for\_invalid\_size()} & silently accepting a size that does
not exist \\
\end{tabularx}

\vspace{6pt}
\noindent Run them with \pkg{pytest}. If they pass, the numbers in the README
are true.

%==============================================================
\section{The on-ramps --- \pkg{examples/}}

Three scripts, in increasing order of ambition.

\begin{tabularx}{\textwidth}{@{}l r Y@{}}
\rowcolor{SchGreen}
\whitebf{Script} & \whitebf{Lines} & \whitebf{What it shows} \\
\pkg{basic\_usage.py} & 77 & \textbf{Start here.} Pure Python, no simulator.
Single-point solve, catalogue-preset twin, full validation table printed to the
terminal. \\
\rowcolor{SchGreenLight}
\pkg{isaacsim6\_actuator.py} & 156 & The \textbf{reusable class} ---
\pkg{HarmonicDriveActuator}. Drop it into any Isaac Sim robot script: it reads
joint state, computes the $\eta$-corrected motor torque, applies it, and logs
the loss breakdown. \\
\pkg{isaacsim6\_pendulum.py} & 270 & The \textbf{complete demo.} A 1-DOF
pendulum (2\,kg arm, 1\,m) driven through a real RT1-H-17-100-CS. Builds the
scene programmatically, runs headless or with GUI, writes CSV. The arm settles
45° $\rightarrow$ 0° in about 3\,s at $\eta$ 80--85\,\%. \\
\end{tabularx}

\vspace{6pt}
\noindent The pendulum example is the one that makes the model \emph{credible
to a roboticist} --- it proves the twin runs inside a real physics engine at
real-time step rates, not only in a notebook.

%==============================================================
\section{How a user actually meets the package}

There are three depths, and most people stop at the first.

\begin{Verbatim}[fontsize=\small,xleftmargin=6pt]
DEPTH 1 · "just give me a number"     → solve_point(30.0, 2000)
          one line, no classes          ~30 seconds to first result

DEPTH 2 · "model my specific gearbox" → twin_for_size(17), or build a
          swap in your own parameters   HarmonicDriveTwin with your own
                                        MeshParams / FrictionParams

DEPTH 3 · "put it in my robot sim"    → copy isaacsim6_actuator.py,
          live actuator in Isaac Sim    call .step() each frame
\end{Verbatim}

%==============================================================
\section{Before you release --- three things to look at}

\begin{watchout}[Open decisions, not blockers]
\begin{enumerate}
  \item \textbf{Schaeffler logo files are tracked in a public MIT repository.}
        \pkg{assets/schaeffler\_wordmark\_green.png} and
        \pkg{we\_pioneer\_motion\_claim.png} are corporate trademarks sitting
        under an MIT license header. Consider shipping an unbranded PDF
        publicly and keeping the branded build internal. The same applies to
        the \emph{INTERNAL --- CONFIDENTIAL} footer on these pages.
  \item \textbf{LaTeX build artefacts were tracked.} The \pkg{.aux},
        \pkg{.fls}, \pkg{.fdb\_latexmk}, \pkg{.log}, \pkg{.out}, \pkg{.toc},
        \pkg{.xdv} and \pkg{build.log} files added churn to every commit. Ship
        \pkg{.pdf} and \pkg{.tex}; ignore the rest. \emph{(Now fixed.)}
  \item \textbf{Uncommitted documentation changes.} Commit or revert everything
        under \pkg{docs/user\_guide/} before tagging a release.
\end{enumerate}
\end{watchout}

\noindent Good news: \pkg{build/}, \pkg{*.egg-info/}, \pkg{.pytest\_cache/} and
\pkg{\_\_pycache\_\_/} are already correctly ignored, and the package tree
itself is clean.

%==============================================================
\section{One-page summary}

\begin{tabularx}{\textwidth}{@{}l r Y l@{}}
\rowcolor{SchGreen}
\whitebf{Piece} & \whitebf{Size} & \whitebf{Role} & \whitebf{User opens it?} \\
\pkg{pyproject.toml} & 37 & shipping label & no \\
\rowcolor{SchGreenLight}
\pkg{README.md} & 249 & shop window & \textbf{yes, first} \\
\pkg{\_\_init\_\_.py} & 129 & front door, public API & indirectly \\
\rowcolor{SchGreenLight}
\pkg{twin.py} & 787 & the engine & if curious \\
\pkg{params.py} & 425 & the dials & \textbf{yes, to adapt} \\
\rowcolor{SchGreenLight}
\pkg{catalogue.py} & 250 & reality check & to validate \\
\pkg{flex\_common/} $\times$ 7 & $\approx$1\,300 & physics foundation &
rarely \\
\rowcolor{SchGreenLight}
\pkg{examples/} $\times$ 3 & 503 & on-ramps & \textbf{yes, to start} \\
\pkg{tests/} & 183 & proof & to verify \\
\rowcolor{SchGreenLight}
\pkg{docs/user\_guide.pdf} & 10 pp & the manual & \textbf{yes} \\
\end{tabularx}

\vspace{10pt}
\noindent Dependencies: \textbf{NumPy and SciPy only.} Matplotlib optional, for
plots. Python $\geq$ 3.10. No compiled extensions, no companion repository, no
license server. That is why it installs anywhere in one command.

\end{document}
